% SCRATCH DRAFT — Vol 2 courtroom coda, Ch28 "Relations, Logic, and the Algebra of Records" (task vol2-coda, author: Podo).
% THIS IS NOT THE BOOK. Do not \input. FLOOR BANKED (integrate on Kodo's accept). Fresh-append (no old beat; body ends line 141).
%
% ⚠ FENCES:
%  - BOX/TYPE-IDENTITY = RESERVE (Kodo turn 30): work the ALGEBRA OPERATIONALLY (records combine/compose/invert/
%    decompose/recompute/correct, all COUNT-PRESERVING). NEVER assert two records are "the same" by
%    indistinguishability (no identity=distinguishability). The channel's "distinguishable symbols" = Shannon
%    CAPACITY (count of symbols), NOT metaphysical identity.
%  - Charge re-sounds: the count is the INVARIANT the algebra preserves (ties Ch27's what-survives-relabeling —
%    now what-survives-operation). ±1/spinor OUT. free-will/continuum-floor/fact->verdict-turn/"you"/juror RESERVED.
%  - Ch29 handoff (Computation) = SEED only ("what machine runs on the records / what can be computed and what
%    cannot / where it fails") — do NOT open FORCING / undecidability / Cantor-Godel (Ch29's LAST fence).
%  - Consensus consolidate (agreement reserved Ch30). "verdict"/2nd-person/penalty-"fine" ABSENT.
%
% CONTINUITY: state carries from Ch27. Ch27 closed handing to "the relations the records build among themselves."
% Ch28 IS that.
%
% PROPOSED SEAM in books/experimentation/latex/chapters/28.tex:
%   KEEP the entire body through line 141 ("...the machine the records run on.").
%   APPEND: \bigskip / \begin{center}\rule{0.4\textwidth}{0.4pt}\end{center} / \bigskip / then the coda below. (Fresh append.)
%
% >>>>>>>>>>>>>>>>>>>>>> APPEND BEGINS (fresh \rule + coda after the body) >>>>>>>>>>>>>>>>>>>>>>

\bigskip
\begin{center}
\rule{0.4\textwidth}{0.4pt}
\end{center}
\bigskip

\noindent The records the case has gathered are not a heap of entries; they carry a structure, an algebra of
their own. Beyond what each mark says stands how the marks stand to one another --- which precedes which, which
excludes which, which depends on which --- and these relations are no loose notes but operations as definite as
arithmetic, that combine and invert and compose. Before the end, the case examines how its own records work upon
each other.

And it finds that everything the records do to one another leaves the count alone. Two relations chain through a
shared middle to make a third, and a relation run backward gives its converse --- but neither adds an entry nor
loses one; they only re-relate the marks already there, and undoing two of them in turn undoes the last one
first, the plainest algebra there is. A record's content stands on an axis of its own,
independent of where and when it was taken --- the same reading able to sit at any place and time, as a book's
text is independent of the shelf that holds it, a separate coordinate and not a separate direction. A channel
need not speak in two states only: refine its alphabet, give it more letters, and the same record is re-spelled
in more symbols without a single mark added, its capacity just the number of distinctions it can hold at once.
On a Christmas Eve more than a century ago a man put his own voice and a violin over the radio, where there had
been only the two-state clatter of Morse --- the first channel to carry not dots and dashes but a continuum of
sound resolved into symbols an ear could read, and not one distinction more than the channel could hold. A
sheet of dependent values, when one is changed, does not pass the change down a wire from cell to cell but
settles the whole into agreement at once, by recalculation and not by a travelling signal. And a record can be
built to survive its own wounding: scatter its order and carry a little parity, and a scratch that would erase a
long passage becomes a sprinkle of single losses, each solved back from the parity. In every one of these the
operation is ours to run; the count it works on, and leaves untouched, is the world's.

So the records are closed under their own algebra: whatever the operations do --- relate, decompose, recompute,
repair --- they act on the entries already there and never add to them. The count is the one thing the algebra
cannot move. This is the invariant of the last chapter met again from a new side: there it was what survives a
change of label; here it is what survives a change of operation. Rename the marks and the count holds; chain
them, invert them, re-spell them in another alphabet, recover them across a scratch, and the count holds still.
The algebra rearranges the record without end and adds to it never --- the count the fixed point of all of it.
It is the sturdiest form yet of the book's one fact: the count is the world's, and everything an instrument does
to a record --- naming it, ordering it, relating it, repairing it --- is the hand's, and leaves the count
exactly where it found it.

And the \emph{charge} is exactly that fixed point. The count the case will collect on passes through every
operation the algebra can perform and comes out unchanged: no relating of the marks, no re-lettering of the
channel, no recalculation of the dependents, no repair across a scratch adds a unit to it or takes one away. The
charge is what the algebra preserves --- the invariant beneath every operation as it was the invariant beneath
every relabeling. What comes to be owed is a count no rearrangement of the record can inflate and no damage can
erase; it survives the handling as it survived the naming. And the trace is the record's whole worked-over
structure --- related, decomposed, recomputed, repaired --- carrying that one count intact through all of it.
Charge and trace are the count and the algebra it rides: the fixed thing, and everything the record does around
it. A charge no operation can move is a charge that will stand however the record is worked --- and that, in the
end, is much of what makes it a charge at all.

And the three readings the case will bring are records of the same kind, closed under the same algebra, their
counts preserved through the same operations. To be compared at all they must be brought into one algebra and
worked upon together --- related, aligned, reconciled --- and through all that working each one's count stays
exactly what it was. The three are corrected and ready, their records robust to every operation that will join
them; the number they must meet on is still not in hand. The algebra can carry the three to the very edge of
their meeting and change not one of their counts on the way.

One examination of the foundations is left, and it is the deepest. Having found the records' order, their
invariant, and their algebra, the instrument asks last what machine runs on them --- what can be computed from a
record and what cannot, how far a finite reckoning reaches and where it fails, and whether the machine can be
trusted to run at all. The next chapter turns to computation itself. The case is still open, its records closed
under their algebra and its three readings not yet agreed; the court does not sit until the record is complete.

% <<<<<<<<<<<<<<<<<<<<<< APPEND ENDS <<<<<<<<<<<<<<<<<<<<<<
